\section{Gamma convolution descent of the original arithmetic source}
\label{sec:gamma-convolution-descent}

\textbf{Purpose and result.}
The Toda calculation expresses the original arithmetic control through
source and relation Gram determinants. The present calculation makes their
input explicit at every tensor degree. It keeps the gamma reference measure
already used for the theta source, multiplies by the actual arithmetic
amplitude, and carries both through the literal sum of the spectral variables.
The sum projection comes with its complete relative-fibre complement.
An exact one-factor coefficient transform supplies every finite source and
relation moment; a gamma majorant supplies an explicit tail bound.

The full proof body of \emph{An exact gamma pushforward for the arithmetic
Toda control} (13 September 2026) follows. Its equations are prefixed GD
and its headings are subordinate here. Its inspection statements describe
the source's pinned historical GitHub read, not a later CI or publication
status. The added final subsection proves the direct analytic map from
the coefficient series to the arithmetic moment function used by Toda.
The arithmetic multiplier, its complex phase, the total mass, the original
cyclic relation, and the supported zero all remain in the displayed maps.

\textbf{Degree domains.}
For a nonzero cyclic module of dimension \(q\), the quotient metrics and
representative maps in GD.31--33 are defined for \(N\geq q-1\).
The adjacent-volume ratios and control expression in GD.34--35 require
\(N\geq q\). At \(N=q-1\) the separate first-degree control formula in
the preceding Toda chapter applies; no \(V_{q-2}^{-1}\) is introduced.
For \(h=1\) the finite quotient is zero-dimensional, while the analytic
seed and its mass remain nonzero. No growing-degree vanishing estimate
or conclusion about all zeta zeros is asserted.

\begingroup
\setlength{\emergencystretch}{3em}
\providecommand{\tightlist}{\setlength{\itemsep}{0pt}\setlength{\parskip}{0pt}}
